\chapter{آرایه‌ها و بردارها}%
\label{ch:arrays}

تا اینجا هر متغیر تنها \emph{یک} مقدار نگه می‌داشت. اما گاهی به مجموعه‌ای از
مقدارها نیاز داریم: نمراتِ یک کلاس، دماهای یک هفته، نام‌های دوستان. برای این
کار از \textbf{آرایه} و \textbf{بردار} استفاده می‌کنیم. در این فصل هر دو را
می‌شناسیم و تفاوتشان را درمی‌یابیم.

\section{آرایه: ردیفی از جعبه‌ها با اندازهٔ ثابت}

آرایه را مانندِ ردیفی از جعبه‌های هم‌گونه تصور کنید که شماره‌گذاری شده‌اند. به
این مثالِ واقعی نگاه کنید:

\begin{salamfa}
روال ریشه:
    نمرات : صحیح[5] = [12, 18, 9, 20, 15]
    ناپایا مجموع : صحیح = 0
    هر نمره از نمرات:
        مجموع = مجموع + نمره
    پایان
    سرچاپ "مجموع:", مجموع
    سرچاپ "میانگین:", مجموع / 5
پایان
\end{salamfa}

در خطِ \code{نمرات صحیح[5] = [12, 18, 9, 20, 15]}:
\begin{itemize}
  \item \code{صحیح[5]} یعنی «آرایه‌ای از ۵ عددِ صحیح».
  \item \code{[12, 18, 9, 20, 15]} پنج مقدارِ اولیه است.
\end{itemize}

\subsection{دسترسی به عناصر با اندیس}

هر عنصرِ آرایه یک \textbf{اندیس} (شماره) دارد. نکتهٔ بسیار مهم: شماره‌گذاری از
\emph{صفر} آغاز می‌شود، نه یک! پس در آرایهٔ بالا:

\begin{center}
\code{نمرات[0]}\,=\,۱۲،\quad \code{نمرات[1]}\,=\,۱۸،\quad \ldots،\quad \code{نمرات[4]}\,=\,۱۵
\end{center}

\begin{warn}
چون اندیس از صفر آغاز می‌شود، در آرایه‌ای با ۵ عنصر، اندیس‌های معتبر ۰ تا ۴
هستند. اگر بکوشید به \code{نمرات[5]} دسترسی پیدا کنید، از مرزِ آرایه بیرون
زده‌اید و این یک خطای رایج و خطرناک است. همیشه دقت کنید که اندیس از
\code{اندازه - 1} بیشتر نشود.
\end{warn}

\subsection{طولِ آرایه با len}

برای دانستنِ اندازهٔ یک آرایه، از روالِ \code{len} استفاده می‌کنیم. این مثالِ
واقعی آن را نشان می‌دهد:

\begin{salamfa}
روال ریشه:
    اعداد : صحیح[5] = [2, 4, 6, 8, 10]
    ناپایا مجموع : صحیح = 0
    ناپایا i : صحیح = 0
    تا i < len(اعداد):
        مجموع = مجموع + اعداد[i]
        i = i + 1
    پایان
    سرچاپ مجموع, len(اعداد)
پایان
\end{salamfa}

با نوشتنِ \code{len(اعداد)} به‌جای عددِ ثابتِ ۵، اگر بعداً اندازهٔ آرایه را
عوض کنید، حلقه همچنان درست کار می‌کند. این عادتِ خوبی است.

\section{برش: پنجره‌ای رو به آرایه}

بسیاری وقت‌ها می‌خواهیم تنها با \emph{بخشی} از یک آرایه کار کنیم چهار عنصرِ
میانی، هرچه پس از عنصرِ سوم می‌آید، یا نیمهٔ نخست بدونِ آنکه آن را در آرایه‌ای
تازه کپی کنیم. \textbf{برش} (slice) یک \emph{نما}ی سبک رو به بازه‌ای از آرایه
است. برش هیچ داده‌ای را در مالکیتِ خود ندارد و هیچ حافظه‌ای روی هیپ نمی‌گیرد؛ در
پشتِ پرده چیزی نیست جز یک اشاره‌گر به نخستین عنصر، به‌همراهِ یک طول. زبان‌هایی مانندِ
Go و Rust نیز همین ایده را دارند.

برش را با یک \textbf{بازه} درونِ کروشه می‌سازید، یعنی \code{a[lo:hi]}:

\begin{salamfa}
روال ریشه:
    اعداد : صحیح[6] = [10, 20, 30, 40, 50, 60]
    میانه : خودکار = اعداد[1:4]
    سرچاپ میانه[0], میانه[1], میانه[2]
    سرچاپ len(میانه)
پایان
\end{salamfa}

\begin{progout}
20 30 40
3
\end{progout}

این بازه \textbf{نیم‌باز} است: \code{اعداد[1:4]} از اندیسِ ۱ \emph{آغاز} می‌شود و
\emph{پیش} از اندیسِ ۴ می‌ایستد، پس اندیس‌های ۱، ۲ و ۳ را دربر می‌گیرد یعنی
طولی برابرِ $4-1=3$. این همان قراردادِ نیم‌بازی است که در شرطِ حلقه
(\code{i < len}) دیدیم و نتیجهٔ تمیزی دارد: طولِ \code{a[lo:hi]} برابرِ
\code{hi - lo} است، و \code{a[lo:hi]} و سپس \code{a[hi:end]} تمامِ بازه را
بی‌شکاف و بی‌هم‌پوشانی پوشش می‌دهند.

\subsection{حذفِ یکی از کرانه‌ها}

اگر کرانهٔ پایین را ننویسید، پیش‌فرض \code{0} است؛ اگر کرانهٔ بالا را ننویسید، تا
انتهای آرایه می‌رود. هر دو را حذف کنید، یعنی \code{a[:]}، و نمایی از کلِ آرایه
می‌گیرید.

\begin{salamfa}
روال ریشه:
    اعداد : صحیح[6] = [10, 20, 30, 40, 50, 60]
    سرچاپ len(اعداد[:3])
    سرچاپ len(اعداد[3:])
    سرچاپ len(اعداد[:])
پایان
\end{salamfa}

\begin{progout}
3
3
6
\end{progout}

\subsection{گونهِ برش، یعنی \texttt{T[:]}}

برش گونهِ خاصِ خود را دارد که \code{T[:]} نوشته می‌شود یعنی «برشی از \code{T}»، در
برابرِ \code{T[N]} که «آرایه‌ای ثابت با \code{N} عنصر» است. این گونه را بیشتر به
صورتِ گونهِ پارامتر می‌بینید، چون روالی که \code{صحیح[:]} می‌گیرد می‌تواند \emph{هر}
رشته‌ای از \code{صحیح}ها را بپذیرد: یک آرایهٔ کامل، یا تنها بخشی از آن.

\begin{salamfa}
روال مجموع(نما : صحیح[:]) : صحیح:
    ناپایا جمع : صحیح = 0
    هر عنصر از نما:
        جمع = جمع + عنصر
    پایان
    برگشت جمع
پایان

روال ریشه:
    اعداد : صحیح[6] = [10, 20, 30, 40, 50, 60]
    سرچاپ مجموع(اعداد[:])
    سرچاپ مجموع(اعداد[2:5])
پایان
\end{salamfa}

\begin{progout}
210
120
\end{progout}

اکنون یک روالِ \code{مجموع} هم روی کلِ آرایه کار می‌کند و هم روی هر برشی از آن.
نه به کپی نیاز دارید و نه باید طول را جداگانه پاس بدهید \code{len} آن را مستقیم
از خودِ برش می‌خواند.

\subsection{برش پنجره است، نه کپی}

از آنجا که برش به‌جای کپی‌کردن، حافظهٔ آرایه را امانت می‌گیرد، نوشتن از طریقِ برش
خودِ آرایهٔ اصلی را تغییر می‌دهد. این تفاوتِ کلیدی با برداشتنِ یک زیرآرایهٔ کپی‌شده
است و دقیقاً همین است که برش را ارزان می‌کند.

\begin{salamfa}
روال ریشه:
    ناپایا اعداد : صحیح[6] = [10, 20, 30, 40, 50, 60]
    ناپایا دنباله : خودکار = اعداد[4:]
    دنباله[0] = 500
    دنباله[1] = 600
    سرچاپ اعداد[4], اعداد[5]
پایان
\end{salamfa}

\begin{progout}
500 600
\end{progout}

مقداردهی به \code{دنباله[0]} از درونِ پنجره دست دراز می‌کند و \code{اعداد[4]} را
بازنویسی می‌کند. برش هیچ داده‌ای از آنِ خود نگه نداشت؛ تنها به آرایه اشاره می‌کرد.

\begin{warn}
برش حافظه‌ای را امانت می‌گیرد که مالکش نیست. تا وقتی برش را به کار می‌برید، آرایهٔ
زیربنایی را زنده نگه دارید، و به یاد داشته باشید که برش طولی \emph{ثابت} دارد:
مانندِ آرایهٔ ثابت رشد نمی‌کند. هر گاه چیزی نیاز داشتید که رشد کند، سراغِ بردار
بروید (در ادامه).
\end{warn}

\section{بردار: آرایه‌ای که رشد می‌کند}

آرایه‌ها یک محدودیت دارند: اندازه‌شان از همان آغاز ثابت است. اما گاهی نمی‌دانیم
چند عنصر خواهیم داشت مثلاً می‌خواهیم تا وقتی کاربر عدد می‌دهد، آن‌ها را نگه
داریم. اینجاست که \textbf{بردار} (Vector) به کار می‌آید: مجموعه‌ای که می‌تواند
\emph{رشد} کند. به این مثالِ واقعی نگاه کنید:

\begin{salamfa}
روال ریشه:
    ناپایا اعداد : وکتور<صحیح> = وکتور {}
    تکرار 1 تا 5 از i:
        اعداد.بیفزا(i * i)
    پایان
    ناپایا مجموع : صحیح = 0
    ناپایا j : صحیح = 0
    تا j < اعداد.طول():
        مجموع = مجموع + اعداد.بگیر(j)[0]
        j = j + 1
    پایان
    سرچاپ "تعداد:", اعداد.طول() برگردان صحیح۶۴
    سرچاپ "مجموع مربع‌ها:", مجموع
پایان
\end{salamfa}

نکته‌های کلیدی:
\begin{itemize}
  \item \code{وکتور<صحیح>} یعنی «برداری از اعدادِ صحیح». بخشِ درونِ
        \code{<\,>} گونهِ عناصر را می‌گوید.
  \item \code{وکتور \{\}} یک بردارِ خالیِ تازه می‌سازد.
  \item \code{.بیفزا(x)} (معادلِ \code{push}) یک عنصرِ تازه به \emph{انتهای}
        بردار می‌افزاید.
  \item \code{.طول()} (معادلِ \code{len}) تعدادِ عناصرِ کنونی را می‌دهد.
  \item \code{.بگیر(j)} (معادلِ \code{get}) عنصرِ شمارهٔ \code{j} را
        برمی‌گرداند. در شیوهٔ کنونیِ زبان، برای خواندنِ مقدار از آن، یک
        \code{[0]} در پی می‌آوریم.
\end{itemize}

\begin{note}
متدها با نقطه (\code{.}) پس از نامِ بردار صدا زده می‌شوند:
\code{اعداد.بیفزا(...)}. این همان شیوه‌ای است که در فصلِ ساختارها (فصلِ ۹) برای
متدهای دلخواهِ خودمان هم به کار می‌بریم.
\end{note}

\subsection{همان بردار به انگلیسی}

نام‌های متدهای بردار مانندِ خودِ گونه‌ها هم فارسی و هم انگلیسی دارند. این
نمونهٔ واقعی همان کار را با نام‌های انگلیسی انجام می‌دهد:

\begin{salamfa}
روال ریشه:
    ناپایا و : Vector<صحیح> = Vector {}
    ناپایا i صحیح = 0
    تا i < 4:
        و.push(i * 10)
        i = i + 1
    پایان
    سرچاپ و.len()
    ناپایا مجموع : صحیح = 0
    i = 0
    تا i < و.len():
        مجموع = مجموع + و.get(i)[0]
        i = i + 1
    پایان
    سرچاپ مجموع
    و.free()
پایان
\end{salamfa}

در پایانِ کار، \code{و.free()} (معادلِ \code{.آزاد()}) حافظهٔ بردار را آزاد
می‌کند. این کار برای بردارهایی که عمرشان به پایان رسیده، یک عادتِ پاکیزه است.

\subsection{بردار به‌مثابهِ پشته}

یک کاربردِ زیبا و رایجِ بردار، ساختنِ \textbf{پشته} (stack) است: ساختاری که
«آخرین چیزی که گذاشتی، نخستین چیزی است که برمی‌داری.» این مثالِ واقعی، با
افزودن و برداشتن، یک پشته می‌سازد:

\begin{salamfa}
روال ریشه:
    ناپایا پشته : وکتور<صحیح> = وکتور {}
    پشته.بیفزا(1)
    پشته.بیفزا(2)
    پشته.بیفزا(3)
    سرچاپ "اندازه =", پشته.طول() برگردان صحیح
    تا پشته.طول() > 0:
        سرچاپ "برداشت", پشته.دربیاور()
    پایان
    سرچاپ "اندازه =", پشته.طول() برگردان صحیح
پایان
\end{salamfa}

اینجا متدِ تازه‌ای می‌بینیم: \code{.دربیاور()} (معادلِ \code{pop}) که آخرین
عنصر را \emph{برمی‌دارد} و برمی‌گرداند. خروجی چنین است: نخست اندازه ۳، سپس
برداشتِ ۳ و ۲ و ۱ (به ترتیبِ معکوس)، و در آخر اندازه ۰.

\section{آرایه یا بردار؟ کدام را برگزینیم؟}

\begin{center}
\begin{tabular}{l l}
\toprule
\textbf{آرایه (\code{T[N]})} & \textbf{بردار (\code{وکتور<T>})} \\
\midrule
اندازهٔ ثابت & اندازهٔ رشدپذیر \\
وقتی تعداد را از پیش می‌دانید & وقتی تعداد نامعلوم است \\
ساده و سریع & انعطاف‌پذیر \\
نیاز به آزادسازیِ دستی ندارد & بهتر است با \code{.آزاد()} آزاد شود \\
\bottomrule
\end{tabular}
\end{center}

\begin{tip}
قاعدهٔ سرانگشتی: اگر تعدادِ عناصر را پیش از اجرا می‌دانید، آرایه برگزینید. اگر
نه مثلاً داده‌ها را یکی‌یکی از کاربر یا از یک فایل می‌گیرید بردار را.
\end{tip}

\begin{deep}[نگاهی به درون: بردار چگونه رشد می‌کند]
بردار درونِ خود یک آرایه نگه می‌دارد. وقتی این آرایه پُر می‌شود و شما عنصرِ
تازه‌ای می‌افزایید، بردار آرایه‌ای \emph{بزرگ‌تر} می‌سازد (معمولاً دوبرابر)،
عناصر را به آن می‌کوچاند و آرایهٔ کهنه را آزاد می‌کند. به همین دلیل افزودنِ یک
عنصر «معمولاً» سریع است، اما گاه‌به‌گاه که آرایه باید بزرگ شود، کمی کندتر
می‌شود. این رفتار که به آن «بزرگ‌سازیِ سرشکن‌شده» می‌گویند در عمل بسیار
کارآمد است.
\end{deep}

\section{خلاصهٔ فصل}

\begin{itemize}
  \item آرایه مجموعه‌ای هم‌گونه با اندازهٔ \emph{ثابت} است: \code{T[N]}.
  \item اندیس‌ها از \emph{صفر} آغاز می‌شوند.
  \item \code{len(آرایه)} طولِ آرایه را می‌دهد.
  \item بردار (\code{وکتور<T>}) مجموعه‌ای \emph{رشدپذیر} است.
  \item متدهای بردار: \code{بیفزا}/\code{push}، \code{بگیر}/\code{get}،
        \code{طول}/\code{len}، \code{دربیاور}/\code{pop}،
        \code{آزاد}/\code{free}.
\end{itemize}

\section{تمرین‌ها}

\exercise آرایه‌ای از ۶ عدد بسازید و بزرگ‌ترینِ آن‌ها را با یک حلقه پیدا کنید.

\exercise آرایه‌ای از دماهای یک هفته بسازید و میانگینشان را حساب کنید.

\exercise برداری بسازید و اعدادِ ۱ تا ۱۰ را در آن بریزید، سپس فقط اعدادِ زوج را
چاپ کنید.

\exercise با کمکِ یک بردار به‌عنوانِ پشته، یک رشته را وارونه کنید (هر نویسه را
در پشته بگذارید و سپس یکی‌یکی بردارید).

\exercise برنامه‌ای بنویسید که آرایه‌ای از اعداد را بگیرد و تعدادِ عناصرِ بزرگ‌تر
از میانگین را بشمارد.
